\documentclass[11pt,a4paper]{article}

\usepackage[T1]{fontenc}
\usepackage[utf8]{inputenc}
\usepackage[ngerman]{babel}
\usepackage{lmodern}
\usepackage{microtype}
\usepackage[a4paper,margin=27mm,headheight=14pt]{geometry}
\usepackage{amsmath,amssymb}
\usepackage{booktabs,tabularx,array}
\usepackage{enumitem}
\usepackage{xcolor}
\usepackage{fancyhdr}
\usepackage{hyperref}

\definecolor{dunkelblau}{HTML}{29495C}
\definecolor{hellblau}{HTML}{EAF1F4}
\definecolor{grau}{HTML}{5B6469}
\definecolor{hellgrau}{HTML}{F3F4F4}

\hypersetup{
  colorlinks=true,
  linkcolor=dunkelblau,
  urlcolor=dunkelblau,
  pdftitle={Vom Fallen zum Grenzwert},
  pdfsubject={Ein reflektierter KI-Einstieg in einen verständnisorientierten Analysisunterricht}
}

\pagestyle{fancy}
\fancyhf{}
\fancyhead[L]{\small\color{grau}Vom Fallen zum Grenzwert}
\fancyhead[R]{\small\color{grau}Reflektierte Arbeit mit KI}
\fancyfoot[C]{\small\color{grau}\thepage}

\setlength{\parindent}{0pt}
\setlength{\parskip}{0.65em}
\setlength{\emergencystretch}{2em}
\raggedbottom

\setlist[itemize]{leftmargin=1.5em,itemsep=0.25em,topsep=0.3em}
\setlist[enumerate]{leftmargin=1.7em,itemsep=0.35em,topsep=0.3em}

\newcommand{\schritt}[1]{\textcolor{dunkelblau}{\textbf{#1}}}
\newenvironment{leitgedanke}
  {\begin{quote}\color{dunkelblau}\itshape}
  {\end{quote}}

\begin{document}

\hypersetup{pageanchor=false}
\begin{titlepage}
  \thispagestyle{empty}
  \vspace*{28mm}

  {\color{dunkelblau}\rule{\textwidth}{1.1pt}}

  \vspace{12mm}

  {\huge\bfseries Vom Fallen zum Grenzwert\par}
  \vspace{5mm}
  {\Large Ein reflektierter KI-Einstieg in einen\\
  verständnisorientierten Analysisunterricht\par}

  \vspace{14mm}

  {\large Das Beispiel des freien Falls\par}

  \vfill

  \begin{minipage}{0.86\textwidth}
    \color{grau}
    Dieses Papier beschreibt einen möglichen Weg. Es ist weder Stundenrezept
    noch Musterlösung. Im Mittelpunkt steht eine didaktische Idee, die von einer
    Lehrkraft entwickelt, im Gespräch mit einer KI geprüft und anschließend zu
    einem Unterrichtsvorschlag ausgearbeitet wird.
  \end{minipage}

  \vspace{15mm}
  {\color{dunkelblau}\rule{\textwidth}{1.1pt}}
\end{titlepage}

\hypersetup{pageanchor=true}
\pagenumbering{roman}
\tableofcontents
\newpage
\pagenumbering{arabic}

\section{Worum es in diesem Beispiel geht}

Der Ausgangspunkt ist unscheinbar: Eine Lehrkraft hält eine Kugel in der Hand
und fragt:

\begin{leitgedanke}
Wie schnell ist die Kugel nach zwei Sekunden, wenn ich sie jetzt fallen lasse?
\end{leitgedanke}

Die Frage ist anschaulich, aber nicht unmittelbar beantwortbar. Wege und Zeiten
lassen sich messen. Aus ihnen können mittlere Geschwindigkeiten berechnet
werden. Die Geschwindigkeit \emph{genau} zu einem Zeitpunkt ist dagegen nicht
einfach ein weiterer Messwert derselben Art.

Damit entsteht ein fachlich ergiebiger Weg:

\[
\boxed{\text{Frage}}
\longrightarrow
\boxed{\text{Daten}}
\longrightarrow
\boxed{\text{Modell}}
\longrightarrow
\boxed{\text{Grenzprozess}}
\longrightarrow
\boxed{\text{Begriff}}
\]

Die Ableitung steht nicht am Anfang. Sie wird benötigt, weil eine verständliche
Frage mit den zunächst verfügbaren Mitteln nicht befriedigend beantwortet werden
kann. In diesem Sinn kann das Beispiel einen genetischen Zugang zur
Momentangeschwindigkeit und damit zum Ableitungsbegriff eröffnen.

Zugleich zeigt die Entstehung dieses Vorschlags eine mögliche Form reflektierter
KI-Nutzung. Die Lehrkraft gibt nicht die Verantwortung für Unterricht an ein
System ab. Sie bringt eine didaktische Entscheidung zur Sprache, lässt sie
prüfen, reagiert auf Einwände, präzisiert ihre Absicht und nutzt die KI
anschließend zur Ausarbeitung.

\begin{leitgedanke}
Die Lehrkraft bleibt Autorin der didaktischen Entscheidung. Die KI kann ihr
Gegenüber sein; sie ersetzt weder Urteil noch Verantwortung.
\end{leitgedanke}


\section{Eine erste didaktische Rückmeldung}

Der Vorschlag besitzt mehrere Stärken.

\subsection{Eine Frage macht Mathematik notwendig}

Die Lernenden beginnen nicht mit der Aufforderung, eine Funktion abzuleiten.
Sie beginnen mit einer Frage, zu deren Beantwortung Mathematik gebraucht wird.
Damit erhält der spätere Begriff eine Funktion: Die momentane Änderungsrate
beantwortet etwas, das mit mittleren Änderungsraten allein offenbleibt.

\subsection{Daten stehen vor der Formel}

Zunächst ist zu klären, welche Größen überhaupt beobachtet und gemessen werden
können. Naheliegend sind Zeitpunkte und die bis dahin zurückgelegten Wege. Aus
ihnen entstehen Tabellen und mittlere Geschwindigkeiten. Das Modell

\[
  s(t)=5t^2
\]

wird erst später eingeführt. Es ist dann nicht bloß eine Rechenvorschrift,
sondern eine Verdichtung und Idealisierung des beobachteten Zusammenhangs.

\subsection{Eine produktive Grenze wird erfahrbar}

Eine endliche Tabelle enthält nur endlich viele Zeitpunkte. Aus je zwei
Messwerten lässt sich eine mittlere Geschwindigkeit für ein Zeitintervall
bestimmen. Kein noch so sorgfältig notierter Tabellenwert ist jedoch schon die
Geschwindigkeit in einem einzelnen Zeitpunkt.

Gerade dieses vorläufige Nicht-Weiterkommen ist didaktisch wertvoll. Es sollte
nicht vorschnell durch die Mitteilung einer Ableitungsregel beendet werden. Die
Grenze des bisherigen Verfahrens erzeugt erst den Bedarf nach einer neuen Idee.

\subsection{Eine reale Stolperstelle gehört dazu}

Nach dem vereinfachten Modell legt die Kugel in zwei Sekunden etwa

\[
  s(2)=5\cdot 2^2=20\,\mathrm{m}
\]

zurück. Das lässt sich in einem gewöhnlichen Klassenraum nicht vorführen. Diese
Beobachtung ist kein nebensächlicher Fehler, sondern eine hilfreiche
Prüffrage. Die Szene mit der Kugel kann als Gedankenexperiment dienen. Für eine
reale Messung bieten sich ein deutlich kürzerer Zeitraum, eine Videoaufnahme
oder vorhandene Messdaten an. Auch hier hilft die Rückfrage eines Gegenübers,
eine gute Idee unter realen Bedingungen tragfähig zu machen.


\section{Ein möglicher Unterrichtsgang}

Die folgenden Phasen bilden keinen verbindlichen Ablauf. Sie markieren
didaktische Funktionen, die je nach Lerngruppe anders angeordnet, verkürzt oder
vertieft werden können.

\subsection{Phase 1: Die Frage öffnen}

Die Lehrkraft zeigt die Kugel und stellt die Ausgangsfrage. Zunächst werden
Vermutungen gesammelt, ohne sofort eine Formel bereitzustellen.

Mögliche Gesprächsimpulse sind:

\begin{itemize}
  \item Was genau bedeutet hier \glqq wie schnell\grqq?
  \item Welche Größen müssten wir kennen oder messen?
  \item Können wir die Geschwindigkeit nach zwei Sekunden direkt messen?
  \item Woran würden wir erkennen, dass eine Antwort gut ist?
\end{itemize}

Bereits hier kann zwischen einer Anzeige auf einem Messgerät und der
mathematischen Bedeutung einer Momentangeschwindigkeit unterschieden werden.
Auch eine Geräteanzeige beruht auf einem Messverfahren, einem endlichen
Zeitfenster und einer Auswertung.

\subsection{Phase 2: Daten erheben oder erhalten}

Die Lernenden erhalten Zeit-Weg-Daten in unterschiedlicher Auflösung. Die Werte
können aus einem geeigneten Versuch, einer Videoanalyse oder einer bewusst
rekonstruierten Messreihe stammen. Im letzten Fall sollte offengelegt werden,
dass es sich nicht um tatsächlich gemessene Daten handelt.

Eine kleine, absichtlich nicht vollkommen modellkonforme Messreihe könnte etwa
so aussehen:

\begin{center}
\begin{tabular}{@{}lrrrrrr@{}}
\toprule
$t$ in $\mathrm{s}$ & 0{,}0 & 0{,}5 & 1{,}0 & 1{,}5 & 1{,}9 & 2{,}0 \\
\midrule
$s$ in $\mathrm{m}$ & 0{,}0 & 1{,}3 & 5{,}1 & 11{,}2 & 18{,}0 & 20{,}1 \\
\bottomrule
\end{tabular}
\end{center}

Die kleinen Abweichungen sind sachlich sinnvoll: Reale Daten tragen
Messunsicherheiten. Sie dürfen nicht nachträglich so behandelt werden, als seien
sie bereits eine exakte Funktion.

\subsection{Phase 3: Mittlere Geschwindigkeiten untersuchen}

Aus zwei Wegmessungen lässt sich die mittlere Geschwindigkeit berechnen:

\[
  \bar v=\frac{\Delta s}{\Delta t}
  =\frac{s(t_2)-s(t_1)}{t_2-t_1}.
\]

Hier liegt eine wichtige Präzisierung. Der Quotient

\[
  \frac{s(t)}{t}
\]

beschreibt bei $s(0)=0$ die mittlere Geschwindigkeit vom Loslassen bis zum
Zeitpunkt $t$. Für die Geschwindigkeit \emph{bei} $t=2$ müssen dagegen
Zeitintervalle betrachtet werden, deren Endpunkte immer näher an $2$ liegen,
zum Beispiel

\[
  [1,2],\qquad [1{,}5,2],\qquad [1{,}9,2].
\]

Die Lernenden erhalten auf diese Weise immer lokalere mittlere
Geschwindigkeiten. Doch mit Messdaten tritt ein neues Problem auf: Je kleiner
das Intervall wird, desto stärker können sich Messfehler im Quotienten
auswirken. Eine feinere Tabelle ist deshalb nicht automatisch eine exaktere
Antwort.

\subsection{Phase 4: Die Grenze der Daten festhalten}

An dieser Stelle lohnt eine bewusste Zwischenbilanz:

\begin{itemize}
  \item Die Daten ermöglichen Aussagen über endliche Zeitintervalle.
  \item Kürzere Intervalle führen näher an den gemeinten Zeitpunkt heran.
  \item Ein einzelner Zeitpunkt besitzt aber keine von null verschiedene
        Zeitdauer.
  \item Endlich viele, ungenaue Messwerte liefern daher keine exakte
        Momentangeschwindigkeit.
\end{itemize}

Die Klasse ist nicht gescheitert. Sie hat die Leistungsfähigkeit und die Grenze
eines Verfahrens erkannt. Nun kann die Frage entstehen, was zusätzlich benötigt
wird.

\subsection{Phase 5: Ein funktionales Modell einführen}

Für den idealisierten freien Fall ohne Luftwiderstand wird das vereinfachte
Modell

\[
  s(t)=5t^2,\qquad t\geq 0,
\]

verwendet. Der Faktor $5$ steht näherungsweise für $g/2$ mit
$g\approx 9{,}81\,\mathrm{m/s^2}$. Die Vereinfachung hält die Rechnung
übersichtlich, ohne die Struktur des Zusammenhangs zu verändern.

Mit der Funktion verändert sich die Art des Gegenstands. Statt einer endlichen
Liste von Messwerten liegt nun ein idealisierter Zusammenhang für alle reellen
Zeitpunkte des betrachteten Intervalls vor. Dadurch können Zeitabstände $h$
beliebig klein gedacht werden.

Für Intervalle, die bei $t=2$ enden, ergibt sich

\begin{align*}
  \frac{s(2)-s(2-h)}{h}
  &=\frac{20-5(2-h)^2}{h}\\
  &=\frac{20-(20-20h+5h^2)}{h}\\
  &=20-5h.
\end{align*}

Die Annäherung wird in einer Modellrechnung sichtbar:

\begin{center}
\begin{tabular}{@{}lrrrr@{}}
\toprule
$h$ in $\mathrm{s}$ & 1 & 0{,}5 & 0{,}1 & 0{,}01 \\
\midrule
$\dfrac{s(2)-s(2-h)}{h}$ in $\mathrm{m/s}$
  & 15 & 17{,}5 & 19{,}5 & 19{,}95 \\
\bottomrule
\end{tabular}
\end{center}

\subsection{Phase 6: Den Grenzprozess vollziehen}

Nun wird nicht einfach ein immer kleineres, aber festes Intervall gewählt. Es
wird gefragt, welchem Wert die mittleren Geschwindigkeiten zustreben, wenn der
Zeitabstand ohne positive Untergrenze verkleinert wird:

\[
  v(2)
  =\lim_{h\to 0}\frac{s(2)-s(2-h)}{h}
  =\lim_{h\to 0}(20-5h)
  =20\,\mathrm{m/s}.
\]

Mit dem üblichen beidseitigen Differenzenquotienten lautet dieselbe Idee:

\[
  \lim_{h\to 0}\frac{s(2+h)-s(2)}{h}
  =\lim_{h\to 0}(20+5h)
  =20\,\mathrm{m/s}.
\]

Der Grenzwert ist nicht der Quotient für $h=0$; dieser wäre gar nicht definiert.
Er beschreibt das Verhalten der Quotienten für beliebig kleine, von null
verschiedene Werte von $h$.

\subsection{Phase 7: Den Begriff sichern}

Erst jetzt wird der fachliche Begriff ausdrücklich gefasst:

\begin{leitgedanke}
Die Momentangeschwindigkeit ist die momentane Änderungsrate der Wegfunktion. Im
mathematischen Modell wird sie als Grenzwert mittlerer Änderungsraten definiert.
\end{leitgedanke}

Für die Wegfunktion $s$ ist die Geschwindigkeitsfunktion ihre Ableitung:

\[
  v(t)=s'(t).
\]

Im vorliegenden Modell gilt $s'(t)=10t$ und damit $v(2)=20\,\mathrm{m/s}$.
Die bekannte Ableitungsregel erscheint hier als ökonomische Fortsetzung einer
inhaltlich bereits verstandenen Grenzwertidee, nicht als deren Ersatz.


\section{Warum die reellen Zahlen eine Rolle spielen}

Messdaten bilden eine endliche Menge. Auch eine sehr umfangreiche Tabelle lässt
zwischen zwei aufeinanderfolgenden Einträgen Lücken. Der Grenzprozess verlangt
dagegen, dass zu jedem noch so kleinen Zeitabstand ein kleinerer betrachtet
werden kann.

Im schulischen Modell wird deshalb von einer Funktion auf einem reellen
Intervall ausgegangen. Das Kontinuum der reellen Zahlen trägt die Idealisierung,
in der $h$ beliebig nahe an null heranrücken kann. Erst dieser Wechsel

\[
  \text{endliche Daten}
  \quad\longrightarrow\quad
  \text{Funktion auf einem reellen Intervall}
\]

ermöglicht den zugehörigen Grenzprozess in der vertrauten Form.

Dabei ist eine fachliche Nuance hilfreich: Nicht die bloße Verwendung des
Symbols $\mathbb{R}$ erzeugt bereits Verständnis. Entscheidend ist die neue
Qualität der Modellierung: Aus einzelnen Beobachtungen wird ein kontinuierlich
gedachter Zusammenhang. Fachmathematisch lassen sich Grenzwerte auch auf
geeigneten dichten Teilmengen untersuchen. Für den Unterricht bleibt der
Übergang von endlich vielen Daten zu einem reellen Funktionsmodell die zentrale
Idee.


\section{Momentangeschwindigkeit als Modellgröße}

Die Rechnung liefert eine präzise Zahl. Trotzdem ist diese Zahl nicht einfach
\glqq die Realität\grqq. Sie gehört zu einem Modell, das bestimmte Aspekte
auswählt und andere ausblendet.

\begin{center}
\begin{tabularx}{\textwidth}{@{}>{\raggedright\arraybackslash}X
                                  >{\raggedright\arraybackslash}X@{}}
\toprule
\textbf{Reale Situation} & \textbf{Mathematisches Modell} \\
\midrule
Kugel mit Ausdehnung, Material und möglicher Rotation
  & idealisierter Körper beziehungsweise Massenpunkt \\
Luftwiderstand und wechselnde Umgebungsbedingungen
  & Luftwiderstand wird vernachlässigt \\
Beschleunigung ist ortsabhängig und wird nur näherungsweise erfasst
  & konstante Beschleunigung, hier vereinfacht mit $g=10\,\mathrm{m/s^2}$ \\
endlich viele Messungen mit begrenzter Genauigkeit
  & Funktionswerte für alle reellen Zeitpunkte eines Intervalls \\
mittlere Geschwindigkeit in einem endlichen Messfenster
  & Momentangeschwindigkeit als Grenzwert \\
\bottomrule
\end{tabularx}
\end{center}

Die Momentangeschwindigkeit ist damit eine mathematisch erzeugte und physikalisch
gedeutete Größe. Sie ist nicht beliebig: Das Modell muss sich an Beobachtungen
bewähren. Sie ist aber auch nicht mit einem unmittelbar gegebenen Stück Realität
gleichzusetzen.

Diese Trennung ist für einen verständnisorientierten Analysisunterricht
wesentlich. Der Grenzprozess macht nicht nur eine Rechnung genauer. Er
konstituiert innerhalb des Modells einen neuen Begriff.


\section{Die Rolle der Lehrkraft und die Rolle der KI}

Das Beispiel ist auch deshalb aufschlussreich, weil der Arbeitsprozess selbst
zum Gegenstand werden kann.

\subsection{Die Lehrkraft entscheidet didaktisch}

Bei der Lehrkraft verbleiben insbesondere die Entscheidungen,

\begin{itemize}
  \item welche Frage für die Lerngruppe bedeutsam und zugänglich ist,
  \item welche fachliche Einsicht angebahnt werden soll,
  \item wann Daten, Darstellungen, Modell und Begriff ins Spiel kommen,
  \item welche Irritation produktiv ist und wo Unterstützung nötig wird,
  \item welche Vereinfachungen verantwortbar sind,
  \item wie Beiträge der Lernenden aufgenommen werden und
  \item woran sich Verstehen zeigen kann.
\end{itemize}

Diese Entscheidungen sind nicht nachträgliche Dekoration eines von der KI
erzeugten Materials. Sie bilden dessen Ausgangspunkt und Maßstab.

\subsection{Die KI kann ein fachliches Gegenüber sein}

Eine KI kann den Entwurf der Lehrkraft unterstützen, indem sie beispielsweise

\begin{itemize}
  \item die innere Linie einer Idee sichtbar macht,
  \item zwischen mittlerer und momentaner Geschwindigkeit unterscheidet,
  \item Rechnungen und Einheiten prüft,
  \item auf die unrealistische Fallhöhe nach zwei Sekunden hinweist,
  \item mögliche Lernschwierigkeiten oder Alternativen benennt,
  \item Daten, Aufgabenstellungen oder Darstellungen variiert und
  \item aus einer geklärten Konzeption ein weitergebbares Dokument erstellt.
\end{itemize}

Auch diese Beiträge müssen geprüft werden. Die KI kennt die konkrete Lerngruppe
nicht, kann didaktische Absichten missverstehen und fachlich plausible, aber für
die Situation unpassende Vorschläge machen.

\subsection{Keine Einbahnstraße}

Reflektierte Arbeit mit KI ist hier ein dialogischer Prozess:

\begin{enumerate}
  \item Die Lehrkraft formuliert ihre Idee und ihre Absicht.
  \item Die KI gibt eine fachliche und didaktische Rückmeldung.
  \item Die Lehrkraft gewichtet die Rückmeldung, widerspricht oder präzisiert.
  \item Die KI arbeitet auf der geklärten Grundlage weiter.
  \item Die Lehrkraft prüft, verändert und verantwortet das Ergebnis.
\end{enumerate}

Die Qualität liegt nicht allein im fertigen Text. Sie liegt auch in den
Rückfragen, Einwänden und Präzisierungen, durch die der didaktische Gedanke
schärfer wird.


\section{Ein möglicher Arbeitsauftrag an eine KI}

Eine Lehrkraft könnte ihren Entwurf zum Beispiel mit einer Anfrage wie der
folgenden zur Diskussion stellen:

\begin{quote}
\small
Ich plane einen Einstieg in die Differentialrechnung über die Frage: \glqq Wie
schnell ist eine fallende Kugel nach zwei Sekunden?\grqq{} Die Lernenden sollen
zunächst überlegen, welche Daten sie benötigen, und aus Zeit-Weg-Daten mittlere
Geschwindigkeiten bestimmen. Erst wenn die Grenze endlicher Messintervalle
sichtbar geworden ist, möchte ich das Modell $s(t)=5t^2$ einführen und zur
Momentangeschwindigkeit als Grenzwert übergehen.

Gib mir bitte zunächst eine kritische didaktische Rückmeldung. Prüfe insbesondere
die Unterscheidung von Messung und Modell, von mittlerer und momentaner
Geschwindigkeit sowie die praktische Plausibilität des Settings. Erstelle noch
kein fertiges Material. Stelle heraus, an welchen Stellen die Lehrkraft eine
didaktische Entscheidung treffen muss. Erst nach meiner Rückmeldung soll daraus
ein Unterrichtsvorschlag entstehen.
\end{quote}

Die Formulierung ist nicht als optimale Prompt-Schablone gemeint. Sie zeigt
vielmehr eine Haltung: Die Lehrkraft legt ihre Überlegungen offen, fordert
Kritik an und behält die Reihenfolge der Arbeit in der Hand.


\section{Mögliche Beobachtungs- und Reflexionsfragen}

Wenn der Unterricht erprobt wird, könnten unter anderem folgende Fragen die
Auswertung leiten:

\begin{itemize}
  \item Unterscheiden die Lernenden die Geschwindigkeit bis zu einem Zeitpunkt
        von der Geschwindigkeit in diesem Zeitpunkt?
  \item Erkennen sie, warum kleinere Messintervalle allein das Problem nicht
        vollständig lösen?
  \item Erleben sie die Funktion als benötigtes Modell oder nur als nachträglich
        mitgeteilte Formel?
  \item Können sie erläutern, was bei $h\to 0$ geschieht, ohne $h=0$ zu setzen?
  \item Trennen sie Messwert, Modellwert und Grenzwert sprachlich und sachlich?
  \item Können sie angeben, welche Annahmen in $s(t)=5t^2$ stecken?
  \item Wird die Ableitung als Antwort auf eine entstandene Frage verstanden?
\end{itemize}

Solche Beobachtungen können wiederum in ein weiteres Gespräch mit der KI
einfließen. Auch dann bleibt die Interpretation der Unterrichtssituation eine
professionelle Aufgabe der Lehrkraft.


\section{Schluss: ein Weg, nicht der einzige}

Das Beispiel des freien Falls ist fachlich elementar und didaktisch keineswegs
belanglos. Es verbindet eine anschauliche Frage mit der Einsicht, dass Daten
allein keine Momentangeschwindigkeit hervorbringen. Erst das funktionale Modell
auf einem reellen Intervall eröffnet den Grenzprozess; in diesem Prozess entsteht
die Momentangeschwindigkeit als Modellgröße.

Ebenso wichtig ist die zweite Ebene des Beispiels. Eine Lehrkraft kann eine
eigene didaktische Idee in das Gespräch mit einer KI geben, ohne ihre
Autorschaft abzugeben. Die KI kann prüfen, irritieren, ordnen und ausarbeiten.
Die Lehrkraft bestimmt jedoch, welche Mathematik aus welchem Grund zum Thema
wird und was für ihre Lernenden ein sinnvoller Weg sein kann.

\begin{leitgedanke}
Nicht die KI entwirft den verständnisorientierten Unterricht. Eine reflektiert
arbeitende Lehrkraft nutzt die KI, um ihre didaktischen Entscheidungen genauer
zu sehen und besser auszuarbeiten.
\end{leitgedanke}

Dieser Weg kann tragen. Er ist nicht der einzige.


\clearpage
\appendix
\section{Die Arbeit geht weiter: didaktische Übergänge gestalten}

Mit dem beschriebenen Unterrichtsgang ist die didaktische Arbeit nicht
abgeschlossen. Im Gegenteil: Gerade an den Übergängen von der Tabelle zum Term,
vom Term zum Grenzprozess und vom Grenzprozess zum Begriff entstehen weitere
Entscheidungen. Sie lassen sich nicht durch ein fertiges Material vorwegnehmen.

Eine Lehrkraft muss für ihre Lerngruppe unter anderem klären:

\begin{itemize}
  \item Wie viele Annäherungsschritte sind hilfreich?
  \item Wann trägt eine weitere Tabelle zum Verständnis bei und wann wird sie
        nur noch mechanisch fortgesetzt?
  \item Wie wird die neue Schreibweise sprachlich und anschaulich vorbereitet?
  \item Welche Vorstellung vom Grenzwert kann zu Beginn aufgebaut werden?
  \item Wie lässt sich erklären, dass $h$ gegen null geht, aber im
        Differenzenquotienten nicht null sein darf?
\end{itemize}

Auch hierbei kann eine KI als Gesprächspartnerin dienen. Die Lehrkraft kann
Varianten prüfen, Formulierungen vergleichen oder mögliche Schwierigkeiten
antizipieren lassen. Sie kann dieselben Entscheidungen aber ebenso im Austausch
mit Kolleginnen und Kollegen oder aus der Beobachtung ihrer Lerngruppe heraus
entwickeln. Entscheidend ist nicht, dass KI an jedem Schritt beteiligt ist.
Entscheidend ist, dass die didaktische Arbeit weitergeht.

\begin{leitgedanke}
Mit oder ohne KI: Die Lehrkraft beobachtet, entscheidet, erprobt, verwirft und
entwickelt weiter. Kein erzeugtes Material nimmt ihr diese Arbeit ab.
\end{leitgedanke}


\subsection{Ist eine weitere Tabelle sinnvoll?}

Eine Tabelle mit noch kleineren Schrittweiten kann sehr hilfreich sein. Sie
macht die Annäherung sinnlich und rechnerisch zugänglich. Ihre Funktion sollte
jedoch geklärt sein.

Zunächst können Messdaten zeigen, dass kürzere Intervalle lokalere mittlere
Geschwindigkeiten liefern. Zugleich bleibt sichtbar, dass Messwerte endlich und
fehlerbehaftet sind. Danach kann eine aus dem Modell $s(t)=5t^2$ berechnete
Tabelle eine andere Aufgabe übernehmen:

\begin{center}
\begin{tabular}{@{}lrrrrrr@{}}
\toprule
$h$ in $\mathrm{s}$
  & 1 & 0{,}5 & 0{,}1 & 0{,}01 & 0{,}001 & 0{,}0001 \\
\midrule
$\dfrac{s(2)-s(2-h)}{h}$ in $\mathrm{m/s}$
  & 15 & 17{,}5 & 19{,}5 & 19{,}95 & 19{,}995 & 19{,}9995 \\
\bottomrule
\end{tabular}
\end{center}

Die Fortsetzung ist dann nicht nur \glqq mehr vom Gleichen\grqq. Sie lenkt den
Blick auf ein stabiles Verhalten: Die Werte kommen $20$ immer näher. Die
Lernenden können vermuten, welcher Wert sich als Ziel der Annäherung anbietet.

Dabei bleiben Fragen für die Lehrkraft offen: Sollen die Werte selbst berechnet
werden? Genügt eine vorgegebene Tabelle? Hilft eine Tabellenkalkulation oder
verdeckt sie an dieser Stelle die Struktur? Ist zusätzlich eine graphische
Darstellung der Sekanten sinnvoll? Die Antwort hängt vom bisherigen Unterricht
und von der Lerngruppe ab.


\subsection{Wie kann die Schreibweise
\texorpdfstring{$f(2+h)$}{f(2+h)} vorbereitet werden?}

Die Schreibweise $f(2+h)$ ist für viele Schülerinnen und Schüler weder
selbstverständlich noch bloß eine kleine Notationsänderung. Mehrere neue
Gedanken treffen aufeinander: Ein Zeitpunkt wird verändert, die Veränderung
erhält einen Namen, und dieser veränderte Zeitpunkt wird als Argument einer
Funktion eingesetzt.

Eine sprachliche Brücke kann deshalb vor der Formel stehen:

\begin{quote}
Wir betrachten den Zeitpunkt $2$ und einen zweiten Zeitpunkt, der nur wenig von
ihm entfernt liegt. Den zeitlichen Abstand zwischen beiden nennen wir $h$. Der
zweite Zeitpunkt ist dann $2+h$.
\end{quote}

Eine einfache Zeitachse kann die Konstruktion stützen:

\[
  2
  \quad\xrightarrow{\hspace*{12mm}h\hspace*{12mm}}\quad
  2+h.
\]

Die beiden zugehörigen Wege sind $s(2)$ und $s(2+h)$. Ihre Differenz wird durch
den zeitlichen Abstand geteilt:

\[
  \frac{s(2+h)-s(2)}{h}.
\]

Für den Einstieg kann auch die rückwärts gerichtete Betrachtung zugänglicher
sein. Dann ist $h>0$ die Länge eines Zeitintervalls, das bei $2$ endet:

\[
  2-h
  \quad\xrightarrow{\hspace*{12mm}h\hspace*{12mm}}\quad
  2,
  \qquad
  \frac{s(2)-s(2-h)}{h}.
\]

Diese Variante hält die Vorstellung des positiven Zeitabstands zunächst fest.
Später kann sie mit der vorwärts gerichteten und schließlich mit der allgemeinen
Schreibweise

\[
  \frac{f(x+h)-f(x)}{h}
\]

verbunden werden. Wann dieser Übergang erfolgt, ist eine didaktische
Entscheidung. Wird die allgemeine Form zu früh eingeführt, können neue Symbolik,
neuer Quotient und neue Grenzwertidee einander überlagern.


\subsection{Wie kann der Grenzwert beginnen, wenn das Grenzwertsymbol noch
unbekannt ist?}

Die Grenzwertvorstellung kann aufgebaut werden, bevor das Symbol $\lim$
eingeführt wird. Tabelle, Term und Sprache tragen zunächst den Gedanken.

Aus

\[
  \frac{s(2)-s(2-h)}{h}=20-5h
\]

und den berechneten Werten lässt sich formulieren:

\begin{quote}
Wir können $h$ immer kleiner wählen. Dadurch können wir die mittlere
Geschwindigkeit beliebig nahe an $20\,\mathrm{m/s}$ heranbringen.
\end{quote}

An dieser Stelle lohnt sprachliche Genauigkeit:

\begin{itemize}
  \item Der Zeitpunkt $2-h$ liegt für kleines $h$ beliebig nahe beim Zeitpunkt
        $2$.
  \item Der Differenzenquotient $20-5h$ liegt dadurch beliebig nahe beim Wert
        $20$.
\end{itemize}

Erst wenn diese Bewegung an Beispielen verstanden ist, erhält sie einen Namen
und ein Zeichen:

\[
  \lim_{h\to 0}(20-5h)=20.
\]

Das Symbol fasst dann einen bereits entwickelten Gedanken zusammen. Es steht
nicht am Anfang als zusätzliche Hürde. Wie weit die Formulierung \glqq beliebig
nahe\grqq{} zu diesem Zeitpunkt präzisiert wird, hängt vom angestrebten
Begriffsniveau ab. Für einen ersten Zugang ist eine tragfähige dynamische
Vorstellung wichtiger als eine vorschnelle formale Definition.


\subsection{Warum darf man nicht einfach
\texorpdfstring{$h=0$}{h=0} setzen?}

Der Differenzenquotient vergleicht Funktionswerte an zwei verschiedenen
Zeitpunkten:

\[
  \frac{s(2+h)-s(2)}{h}.
\]

Für $h=0$ fallen die beiden Zeitpunkte zusammen. Es entstünde

\[
  \frac{s(2)-s(2)}{0}=\frac{0}{0}.
\]

Dieser Quotient ist nicht definiert. Zugleich verschwindet das Zeitintervall,
über das eine mittlere Geschwindigkeit bestimmt werden könnte. Deshalb wird
$h=0$ nicht in den Differenzenquotienten eingesetzt.

Im konkreten Beispiel gilt für $h\neq 0$:

\begin{align*}
  \frac{s(2+h)-s(2)}{h}
  &=\frac{5(2+h)^2-20}{h}\\
  &=\frac{20h+5h^2}{h}\\
  &=20+5h.
\end{align*}

Der letzte Term ist auch bei $h=0$ definiert. Daraus kann leicht der Eindruck
entstehen, man habe nun doch einfach $h=0$ eingesetzt. Inhaltlich muss jedoch
unterschieden werden:

\begin{itemize}
  \item Der ursprüngliche Differenzenquotient bleibt bei $h=0$ undefiniert.
  \item Für jedes $h\neq0$ besitzt er denselben Wert wie $20+5h$.
  \item Das Verhalten von $20+5h$ in der Nähe von null zeigt, welchem Wert der
        Differenzenquotient zustrebt.
\end{itemize}

Eine knappe Formulierung für die Lerngruppe kann lauten:

\begin{leitgedanke}
Wir berechnen den Quotienten nicht bei $h=0$. Wir untersuchen, was mit seinen
Werten geschieht, wenn $h$ beliebig nahe an null heranrückt.
\end{leitgedanke}

Der Grenzwert schließt damit nicht einfach eine rechnerische Lücke. Er beschreibt
das Verhalten in der Umgebung einer Stelle, obwohl der untersuchte Quotient an
der Stelle selbst nicht definiert ist.


\subsection{Weitere Entscheidungen, die offenbleiben dürfen}

Auch nach diesen Klärungen ist kein vollständig festgelegter Unterricht
entstanden. Weitere Fragen können produktiv offenbleiben:

\begin{itemize}
  \item Wird zunächst nur ein einseitiges Annähern betrachtet oder werden beide
        Seiten früh miteinander verglichen?
  \item Wann wird die Sekante im Weg-Zeit-Diagramm einbezogen, wann die Tangente?
  \item Wie werden Messunsicherheit und Modellabweichung thematisiert, ohne den
        begrifflichen Aufbau zu überladen?
  \item Wann wird vom konkreten Zeitpunkt $2$ zum allgemeinen Zeitpunkt $t$
        übergegangen?
  \item Wann wird aus der einzelnen Momentangeschwindigkeit eine
        Geschwindigkeitsfunktion?
  \item Welche Äußerungen der Lernenden zeigen tatsächlich, dass der
        Grenzprozess verstanden wurde?
\end{itemize}

Solche Fragen sind kein Mangel des Materials. Sie sind Ausdruck professioneller
Autorschaft. Eine KI kann beim Weiterdenken helfen, Alternativen erzeugen und
Widersprüche sichtbar machen. Sie kann aber weder die Reaktionen einer konkreten
Lerngruppe vorwegnehmen noch die daraus folgenden Entscheidungen verantworten.

\begin{leitgedanke}
Das Gespräch mit der KI kann enden; die didaktische Arbeit der Lehrkraft geht
weiter. Und auch wenn das Gespräch fortgesetzt wird, bleibt es die Lehrkraft,
die beobachtet, urteilt und entscheidet.
\end{leitgedanke}

\end{document}
